\chapter{Programa de inspecciones}
\label{cap:mm-inspecciones}

\opening{El intervalo máximo es de 100 horas o 12 meses, lo que ocurra primero.}

\section{Inspecciones programadas}
\fuente{MM, capítulos 1 y 2}

\begin{table}[htbp]
\centering\scriptsize
\caption{Inspecciones y tareas programadas (MM, tabla 1). Consulte la última revisión de cada publicación para su aplicabilidad. Los intervalos marcados con $\ast$ son más cortos que los que indica el documento referido.}
\label{tab:mm-programadas}
\setlength{\tabcolsep}{2.5pt}
\begin{tabularx}{\linewidth}{@{}Y>{\raggedright\arraybackslash}p{2.7cm}@{}}
\toprule
Tarea & Intervalo \\
\midrule
Inspección según el manual del operador de Lycoming & Primeras 25 h; cada 50 h; cada 100 h; cada 300 h$\ast$ \\
Lycoming SB 480: cambio de aceite y filtro, limpieza de filtros; inspección del contenido del filtro/malla (los filtros en adaptadores D723-1 no requieren alambre de frenado) & Primeras 25 h; cada 50 h; cada 4 meses \\
Lycoming SI 1129: tensión de la correa del alternador & Primeras 25 h; cada 100 h \\
Lycoming SB 301: procedimientos y límites de servicio de válvulas & Primeras 100 h$\ast$; cada 300 h$\ast$ \\
Lycoming SB 388 (también cilindros de reemplazo): estado de válvulas y guías de escape & Primeras 100 h$\ast$; cada 300 h$\ast$ \\
Lycoming SB 366, si aplica: tornillos del cuerpo del acelerador del carburador & Cada 50 h \\
Lycoming SI 1191: compresión de cilindros & Cada 100 h \\
CMI SB 643, si aplica: intervalos de mantenimiento de magnetos CMI/TCM/Bendix & Cada 100 h; cada 500 h; cada 5 años; cada 12 años \\
Inspección de 100 horas/anual (\S\,2.400) & Cada 100 h; cada 12 meses \\
Mantenimiento de las puntas de palas del rotor principal (\S\,26-60) & Cada 100 h; cada 12 meses \\
Lubricar rodamientos A181-4 rev.~K, L y M (\S\,22-42) & Cada 300 h; cada 12 meses \\
Lubricar rodamiento A184 (\S\,22-41) & Cada 300 h; cada 12 meses \\
Limpiar detectores de virutas de las cajas (\S\S\,22-11 y 22-21) & Cada 500 h; cada 12 meses \\
Pesar el extintor B381-1 o -4 (\S\,15-62A) & Cada 12 meses \\
CMI SB 663: platinos del tacómetro de magneto de dos cables P/N 10-400507 & Cada 500 h \\
Inspección y servicio del lubricante del embrague (\S\,22-30) & Cada 500 h \\
Drenar y lavar las cajas de transmisión (\S\S\,22-13 y 22-23) & Cada 500 h \\
FAA AD 88-26-01 R2 (solo espigas A158-1) & Cada 500 h \\
Lubricar rodamientos A181-4 rev.~N y posteriores (\S\,22-42) & Cada 800 h; cada 4 años \\
Inspección de 2200 horas/12 años (\S\,2.600) & Cada 2200 h; cada 12 años \\
\bottomrule
\end{tabularx}
\end{table}

\section{Introducción a la inspección}

El R22 debe inspeccionarse periódicamente para verificar que está aeronavegable. \textbf{El intervalo máximo es de 100 horas de tiempo en servicio o 12 meses calendario (anual), lo que ocurra primero}; puede extenderse hasta 10 horas, sin acumulación, si la reglamentación local lo permite. Las pérdidas de fluidos, decoloraciones, abolladuras, rayas, mellas, grietas, excoriaciones, roces, \emph{fretting} y corrosión justifican una investigación adicional; los elementos no aeronavegables se reemplazan o reparan según lo permita RHC.

\section{Procedimientos generales}

Para la inspección visual, vea la AC 43.13-1B de la FAA (capítulo 5, sección 2), y para el reemplazo de cableado, su capítulo 11, sección 8, párrafo 11-97. Cuando se requiera, la inspección por partículas magnéticas se hace según ASTM E~1444 y MIL-STD-1907, y la de líquidos penetrantes fluorescentes según ASTM E~1417 y MIL-STD-1907.

\subsection{Rodamientos de bolas y de rodillos}

El primer indicio de falla de un rodamiento suele ser un aumento del ruido, que casi siempre comienza varias horas antes de la falla o de cualquier aumento de temperatura. Escuche el sistema de transmisión en el arranque y la detención: un rodamiento en falla produce un fuerte zumbido, retumbo, gruñido o sonido de sirena. Al oír un ruido inusual, inspeccione a fondo todos los rodamientos antes de volar. Un rodamiento en falla puede tener el sello deformado o expulsar mucha grasa. No confíe en los Telatemps: el aumento de temperatura puede ocurrir solo segundos antes de que el rodamiento se desintegre (vea \S\,22-40).

La falla en vuelo de cualquiera de los rodamientos del actuador podría causar la pérdida de potencia al rotor y un accidente grave. El rodamiento superior (de bolas) está en el eje del embrague, detrás de la polea superior; el inferior (de rodillos), en el eje del ventilador, detrás de la polea inferior. Justo antes de su falla completa, la luz del embrague puede parpadear constantemente (encenderse y apagarse en menos de un segundo), lo que no debe confundirse con el retensado normal en vuelo (encendida de 1 a 8 segundos y luego apagada). No reanude el vuelo hasta determinar la causa del parpadeo.

\subsection{Tubos de empuje--tracción}

Las mellas, cortes o rayas de no más de 0,010~in de profundidad y no más de 1/4 de la circunferencia del tubo pueden pulirse longitudinalmente con lija al agua de grano 320 o más fino, con un radio de transición mínimo de 1~in. Si la profundidad excede estos límites, el tubo debe reemplazarse. Si el tubo está abollado o aplastado más del 5\,\% de su diámetro, debe reemplazarse.

\subsection{Terminales de varilla y rodamientos esféricos}

\begin{checklist}
\chk{Juego axial máximo}{0,020 in}
\chk{Juego radial máximo}{0,010 in}
\chk{Holgura entre la pista exterior y la carcasa del terminal}{No permitida}
\chk{Terminales no remachados con agujero testigo}{Deben impedir el paso de un alambre de 0,020 in}
\end{checklist}

Extensión máxima del terminal cuando no hay agujero testigo: A101-4 y F101-4, 1,10~in; D163-1, 1,40~in; D173-1, 1,30~in.

Las contratuercas (\emph{jam nuts}) y tuercas \emph{palnut} de los terminales se aprietan según \S\,23-32 y se marcan con franja de torque en la posición más visible para el prevuelo. La franja debe cruzar las tuercas hasta el vástago del terminal y el tubo (o el cilindro de la varilla de paso, la horquilla, el soporte, el montante, etc.). Las franjas se deterioran y deben renovarse periódicamente. Los terminales deben quedar centrados, es decir, posicionados para permitir el mayor movimiento posible sin trabarse.

\begin{advertencia}
El armado de los controles de vuelo es crítico y requiere inspección por una persona calificada. Si no hay una segunda persona disponible, quien los armó debe tomarse un descanso de 5 minutos antes de inspeccionar las conexiones que montó.
\end{advertencia}

\begin{precaucion}
Los rodamientos con camisa de teflón no deben lubricarse ni limpiarse con solvente.
\end{precaucion}

\subsection{Rodamientos elastoméricos}

Se usan en el cubo A062-1 del rotor de cola. La fatiga, la contaminación con aceite o la sobrecarga pueden degradar el elastómero.

\begin{itemize}
\item \textbf{Fatiga:} las pequeñas grietas superficiales y el polvillo de elastómero (``migas de goma'') son normales y no exigen reemplazo. Al crecer las grietas se pierde elastómero, con menor rigidez y mayor vibración. \textbf{Reemplace el rodamiento si hay grietas profundas (más de 0,10~in) en más del 25\,\% de la cara del elastómero.}
\item \textbf{Contaminación:} evite la exposición a aceite, grasa, fluido hidráulico, solventes de limpieza y antioxidantes; lávelos de inmediato con detergente y agua. Un rodamiento contaminado muestra hinchazón, bordes ondulados o despegue, y debe reemplazarse.
\item \textbf{Sobrecarga:} ocurre al superarse la resistencia a la tracción del elastómero o de su unión con el metal, a menudo por cargas normales sobre un rodamiento debilitado por fatiga o contaminación. Se reconoce por grandes extrusiones y grietas grandes y limpias.
\item \textbf{Despegue:} si el elastómero se separa de los bujes metálicos en más del 25\,\% del área unida, reemplace el rodamiento.
\end{itemize}

\subsection{Indicadores Telatemp}

Los Telatemps autoadhesivos registran los cambios de temperatura de operación de rodamientos y cajas. Trace una línea de referencia entre el cuadro de mayor temperatura que se oscureció en operación normal y el siguiente sin oscurecer. En cada revisión posterior, determine si se ennegreció un cuadro adicional; si el aumento no puede explicarse por un cambio en las condiciones de operación, examine cuidadosamente el componente antes de volar.

\begin{nota}
Los Telatemps contaminados con un derivado del petróleo pueden indicar erróneamente, típicamente con cuadros blancos sin activar entre cuadros oscurecidos a ambos lados. Reemplácelos y limpie el área con acetona antes de instalar el nuevo.
\end{nota}

\begin{table}[htbp]
\centering\small
\caption{Telatemps del R22.}
\begin{tabularx}{\linewidth}{@{}lY@{}}
\toprule
Telatemp (rango) & Ubicación \\
\midrule
F110-4 (104 a 132~°C / 220 a 270~°F) & Caja del rotor principal A006-1 \\
F110-2 (60 a 88~°C / 140 a 190~°F) & Rodamiento superior del actuador A184-3 (detrás de la polea superior) \\
F110-2 & Rodamiento inferior del actuador A181-4 (detrás de la polea inferior) \\
F110-2 & Caja del rotor de cola A021-1 o B021-1 \\
\bottomrule
\end{tabularx}
\end{table}

\subsection{Marcas de torque}

Si la marca de torque que queda en un sujetador no basta para juzgar la integridad de la unión, retire la \emph{palnut} si corresponde y aplique el torque especificado. Si el sujetador se mueve, desarme la unión e inspeccione las piezas (\emph{fretting}, deformación de roscas, agujeros ovalados, etc.) y reemplace las dañadas. Si no se mueve, instale una \emph{palnut} nueva si corresponde, aplique el torque estándar (\S\,23-32) y marque la franja de torque.

\section{Chequeos en tierra y en vuelo de la inspección de 100 horas/anual}

Complete estas listas junto con la inspección de 100 horas o anual, anotando y corrigiendo cualquier discrepancia.

\subsection{Chequeo en tierra}

\begin{pasos}
\item \textbf{Acelerador:} libre giro con el colectivo completamente abajo y completamente arriba.
\item \textbf{Resorte de sobrerrecorrido:} gire el acelerador más allá del ralentí contra el tope; al soltarlo debe volver de inmediato a la posición normal de ralentí.
\item \textbf{Mezcla:} funcionamiento suave, sin trabas; botón de destrabe funcionando; retroceso de 0,03 a 0,10~in en la posición rica a fondo.
\item \textbf{Calefacción del carburador:} funcionamiento suave; retroceso de 0,03 a 0,10~in en la posición completamente cerrada.
\item \textbf{Cíclico:} libre en todo el recorrido con fricción y compensación derecha apagadas. El pomo de fricción debe girar entre 1/8 y 1 vuelta en sentido horario antes de empezar a frenar. Recorrido completo con fricción y compensación aplicadas. \textbf{Fricción máxima de 10~lb} medida lateralmente en la empuñadura, con compensación derecha apagada, fricción aplicada y varillas de paso del rotor principal de proa a popa.
\item \textbf{Colectivo:} libre en todo el recorrido sin fricción; con fricción aplicada, \textbf{12 a 20~lb} medidas en la empuñadura.
\item \textbf{Asistente de calefacción (si está instalado):} con colectivo abajo y calefacción total, suba el colectivo a fondo: la calefacción debe cerrarse. Bájelo a fondo: debe abrirse por completo. Sin fricción de colectivo, cierre la calefacción y verifique que el colectivo permanezca abajo.
\item \textbf{Pedales:} funcionamiento suave, sin trabas.
\item \textbf{Controles removibles:} sujetadores firmes.
\item \textbf{Iluminación e instrumentos} (interruptor \textsf{MASTER} en ON): luz de CO destella dos veces; CAT aproximadamente igual a OAT; luces \textsf{ALT} y \textsf{OIL} encendidas; indicadores de combustible indican; luces de navegación, panel, estroboscópica, de aterrizaje (con embrague acoplado) y de mapa funcionan; amperímetro indica descarga; temperatura de aceite y CHT con leve desviación con motor frío; \textsf{MR CHIP}, \textsf{MR TEMP} y \textsf{TR CHIP} se encienden con el sensor puesto a masa o el botón de prueba presionado; \textsf{LOW FUEL} se enciende tras aproximadamente 1~s al presionar el sensor del tanque o el botón de prueba.
\end{pasos}

\subsection{Prueba de motor}

\begin{pasos}
\item Efectúe el prevuelo del POH.
\item \textsf{MASTER} en ON. Aplique el freno rotor: luz \textsf{BRAKE} encendida; con el freno aplicado, el arranque no debe accionarse al girar la llave a START.
\item Arranque y haga la prueba del POH, pero con las fricciones de cíclico y colectivo aplicadas y el gobernador desconectado. Tras acoplar el embrague, \textbf{las palas deben girar dentro de 5~s} y la luz \textsf{CLUTCH} apagarse \textbf{dentro de 100~s}.
\item Amperímetro indica carga; luz \textsf{ALT} apagada.
\item Ambos magnetos se ponen a masa (momentáneamente en OFF) a 60\,\% de RPM.
\item Con alternador y batería en OFF: los tacómetros funcionan y todas las luces se apagan (diodo de aislamiento de la barra de tacómetros funcionando).
\item Sin ruidos anormales de rodamientos al variar la RPM; el mecánico escucha cerca de la transmisión por correas.
\item A 75\,\%, gobernador ON; suba a 85\,\% y suelte: el gobernador debe llevar la RPM a 104\,\% $\pm$ 1,5\,\%. Suba a 106\,\% y suelte: debe bajarla a 104\,\% $\pm$ 1,5\,\%.
\item Agujas de motor y rotor dentro del 1\,\% entre sí a 104\,\%.
\item Voltaje del alternador: 13,4 a 13,9~V con unidad de control A942-3; 13,0 a 13,7~V con A942-1 o -2.
\item Calefacción funcionando (si está instalada).
\item Las agujas de los tacómetros no saltan más del 2\,\% al transmitir en 118,00, 125,00 y 136,975~MHz con el gobernador ON.
\item Suba el colectivo 0,5~in en la empuñadura y reduzca lentamente la RPM: bocina y luz de baja RPM se activan entre 97 y 96\,\% y siguen activas hasta el ralentí.
\item Ralentí con motor caliente, embrague acoplado y acelerador cerrado: 53 a 57\,\%.
\item Mezcla de ralentí (motor caliente, embrague desacoplado, acelerador cerrado): aumento de 2 a 4\,\% de RPM al cortar lentamente la mezcla. Ajuste el tornillo de mezcla de ralentí; si no logra el aumento, póngalo a 1½ vueltas desde completamente adentro y ajuste para un ralentí suave.
\end{pasos}

\subsection{Chequeo en vuelo}

\begin{pasos}
\item \textbf{Estacionario:} todos los indicadores en verde; controlabilidad en giros a izquierda y derecha; niveles de vibración.
\item \textbf{Vuelo nivelado} (unos 3000~ft de altitud de densidad si es posible, MCP, gobernador ON, compensación derecha ON, fricciones OFF): niveles de vibración; el cíclico no se desplaza longitudinal ni lateralmente; el colectivo no sube ni baja solo; posición de los pedales de 0,0 a 0,5~in a la izquierda.
\end{pasos}

\section{Preparación de la célula}

La célula debe limpiarse a fondo antes de la inspección (14 CFR 43, apéndice D, párrafo a): palas, cubos y exterior con una solución de jabón suave y agua (vea el capítulo~\ref{cap:manejo}). Anote cualquier pérdida de fluido antes de limpiar.

\begin{precaucion}
No rocíe con agua a presión los magnetos, el cubo del rotor principal, el respiradero de la caja de cola ni la zona del plato oscilante: el agua puede entrar y causar corrosión y degradación de los lubricantes.
\end{precaucion}

\section{Inspección de 100 horas/anual de la célula}

RHC recomienda guardar una copia de la última lista completada con los registros de mantenimiento. Los números de zona remiten a las figuras de acceso del MM (figuras 2-3A y 2-3B). En cada zona aplica además el ítem común: \emph{sujetadores y marcas de torque} (verificar la firmeza de todos y renovar las franjas deterioradas) y \emph{cableado} (sin cables sueltos, rotos ni rozados; ordenado y firme).

\subsection*{1. Retén de la luz de aterrizaje}
\begin{itemize}
\item No retire las tapas de los bloques de rodamiento de los pedales salvo que la prueba funcional indique un problema.
\item Retén y soporte de la luz: estado y firmeza.
\item \textbf{Bloques de rodamiento de los pedales} (con luz y espejo): juego máximo \textbf{0,050~in axial y 0,020~in radial}; soportes verticales de chapa sin grietas ni \emph{fretting}.
\item Pedales sin grietas; tubos de empuje--tracción sin grietas en los extremos, terminales según \S\,2.120.
\item Estructura de chapa sin corrosión ni grietas; caja de batería en la nariz (si existe) sin fugas de electrolito; mentón sin equipos sueltos ni residuos.
\end{itemize}

\subsection*{2. Consola superior}
\begin{itemize}
\item No es necesario retirar la consola de instrumentos (\S\,13-70) para las inspecciones programadas: basta retirar el retén de la luz de aterrizaje o la aviónica.
\item Consola, bisagras, instrumentos y conexiones; fuente del estrobo y su cableado (sobre el cortafuego horizontal o detrás del respaldo derecho) sin grietas en el montaje; soporte B255-1 y su pasacables (solo consola de 10 agujeros).
\item Tubos del rotor de cola visibles con holgura de operación; batería y caja (si está en la nariz): nivel y densidad del electrolito (solo baterías no selladas), conexiones y estructura sin corrosión; bandejas de radio firmes.
\end{itemize}

\subsection*{3. Panel de disyuntores}
\begin{advertencia}
Asegúrese de que el sistema eléctrico permanezca apagado con el panel abierto.
\end{advertencia}
\begin{itemize}
\item Cableado y conexiones sin corrosión; disyuntores firmes (revise la aplicabilidad de las AD); \textbf{collar rojo en el botón del disyuntor \textsf{CLUTCH}}.
\item Barras de cobre sin corrosión ni dobleces y sin contacto con la estructura; interior limpio.
\end{itemize}

\subsection*{4. Cubiertas horizontal, de la caja del cíclico, del vientre y vertical}
Para retirar la cubierta de la caja del cíclico: desenrosque el pomo de ajuste de la compensación, retire el pasador elástico del pomo de fricción y, con asistente de calefacción, desenrosque su pomo.
\begin{itemize}
\item Letreros legibles. Caja del cíclico sin grietas, remaches sueltos ni corrosión.
\item Bastón del cíclico sin grietas en las soldaduras; \textbf{desgaste máximo de 0,015~in} en las placas de la base del bastón por la carcasa del terminal delantero del tubo A121-1.
\item Pivote del cíclico (extremo delantero del tubo de torsión) y horquilla (extremo trasero): rodamientos según \S\,2.120.
\item Compensación derecha: con las varillas de paso de proa a popa, al activarla debe producir fuerza a la derecha; con el cíclico a fondo a la izquierda, las espiras del resorte no deben trabarse.
\item Fricción del cíclico, tubos horizontales del cíclico (sin grietas junto a los bloques de refuerzo), tubos del rotor de cola y balancín inferior; cables de mezcla y calefacción firmes; balancín y tubo del asistente de calefacción.
\item Líneas pitot y estática: sin grietas, roces ni dobleces; retire los tapones de drenaje de las tes, drene el agua y reinstálelos.
\item \textbf{Cordón elástico de compensación:} retírelo y estírelo buscando huecos (hebras rotas); verifique los anillos de traba de los ganchos. (No se requiere si el cíclico no necesita compensación hacia adelante en crucero.)
\item Tubos verticales del rotor de cola y del cíclico: se permite el contacto entre los extremos inferiores de los tubos A121-3 y -7 con el cíclico en las esquinas traseras de la caja. Balancín, horquilla del cíclico y resorte de sobrerrecorrido (se permite su contacto con el tubo A121-1 con el cíclico en la esquina delantera derecha).
\item Fricción y tope del colectivo sin mellas, cortes ni rayas, sin flexión ni trabas; funda del colectivo sin roturas.
\end{itemize}

\subsection*{5. Respaldos y protector del resorte del colectivo}
\begin{itemize}
\item Bastidor superior (partes visibles); cableado y antenas con pasacables bien instalados; líneas pitot y estática.
\item Colectivo: rodamientos según \S\,2.120, sin grietas en soldaduras, bastón removible firme y retén A783-1 sujeto con amarras.
\item Varilla de conexión del acelerador sin grietas. \textbf{Gobernador:} terminales según \S\,2.120 y holgura suficiente del cableado con el colectivo arriba.
\item \textbf{Resorte del colectivo:} sin trabas ni asperezas; varillas guía lubricadas; espiras sin tocarse con el colectivo abajo; agujeros de la tapa inferior sin ovalarse.
\item Drenajes del cortafuego despejados y sin abrazaderas que corten los tubos; controlador del gobernador/EMU y su cableado (\S\,33-137).
\item Salida del tanque principal, manguera, válvula de corte (funcionamiento suave) y manguera del auxiliar sin fugas y con holgura respecto de la estructura.
\end{itemize}

\subsection*{6. Puertas del capó y capó trasero}
\begin{itemize}
\item Fusibles del actuador: tipo \textbf{AGC-1½}, sin corrosión, traba de giro funcionando. Antenas, ELT (14 CFR 91.207(d) si corresponde, correa D693-4 abrochada, holgura con la transmisión), línea pitot y toma estática limpia.
\item Tanques: sin fugas; interior sin objetos extraños ni grietas en el deflector. Manguera del auxiliar tipo ``desprendible'' C595-1: acople sin fugas, cuatro remaches firmes y separación máxima de 0,005~in.
\item Drenajes de sumidero: abren fácil, drenan libremente, cierran por resorte y sellan. \textbf{Aviso de bajo combustible:} con la batería ON, presione suavemente el flotador del sensor con una varilla de madera limpia: la luz debe encenderse tras aproximadamente 1~s.
\item \textbf{Tapas de combustible:} con tanques de vejiga, la tapa A689-6 debe ser rev.~J o posterior; empaquetadura sin deterioro; calcomanía legible y marcas de alineación coincidentes al cerrar (R22 SB-107).
\item Soportes y eje intermedio (\emph{jackshaft}) del colectivo sin grietas ni corrosión (si es abierto, sin corrosión interna; extremos de los montantes sellados). Tubos del plato oscilante: sin roturas ni desgaste de las fundas termorretráctiles de los tubos delanteros. Guías o rodillos limpios, sin desgaste del metal y con rodillos libres.
\item Tubo del rotor de cola y balancín delantero: sin grietas en el montaje de chapa del balancín. Ducto y tobera de enfriamiento de la caja principal sin roturas ni obstrucciones; pasacables del cortafuego horizontal.
\item \textbf{Caja principal:} soportes, unión del tubo del mástil estático y fundición inferior; Telatemp; nivel de aceite sin fugas; emisores de efecto Hall firmes. Cada 500~h o anualmente, retire y limpie el detector de virutas (\S\,22-11); drene y lave la caja a intervalos que no excedan 500~h (\S\,22-13).
\item \textbf{Freno rotor:} cable, polea y microinterruptor; \textbf{espesor mínimo de las pastillas 0,030~in}; holgura con la horquilla de entrada con el freno suelto.
\item Horquilla de entrada de la caja principal con sus imanes firmes.
\item \textbf{Placa flexible delantera:} sin deformaciones, mellas, rayas, grietas, corrosión ni \emph{fretting} (si hay \emph{fretting}, consulte a RHC); arandelas pegadas a ambos lados de cada oreja. La A947-1 debe ser rev.~E o posterior.
\item Horquilla delantera del eje del embrague sin grietas, corrosión ni \emph{fretting}. \textbf{Embrague:} sin fugas de aceite; sin corrosión del eje, sobre todo en las uniones con los sellos A966-1 (elimine la corrosión superficial y aplique cera o compuesto anticorrosivo); placas retenedoras de sellos sin grietas. Se recomienda inspeccionar y dar servicio al lubricante cada 500~h.
\item \textbf{Polea superior:} superficies de contacto lisas. Reemplácela si tiene corrosión, picaduras, descascarado, asperezas, crestas filosas, desgaste a través del anodizado o ampollas del recubrimiento metalizado opcional.
\item \textbf{Bastidor superior de tubos de acero:} sin grietas ni corrosión (si hay corrosión, \S\,2.560); inspeccione cada soldadura con luz y espejo. Está cargado a fatiga: revise todas las uniones a fondo.
\item Cortafuego horizontal sin grietas, pandeos ni arrugas; mamparo trasero de la cabina sin corrosión, remaches sueltos, deformación ni grietas.
\end{itemize}

\begin{advertencia}
Las placas flexibles A193, sin arandelas pegadas, son obsoletas y deben reemplazarse por placas A947 con arandelas pegadas. Si una arandela pegada se separa de una placa A947, la placa queda no aeronavegable y no puede repararse.
\end{advertencia}

\subsection*{7. Faldones laterales, tapas de bujías y caja de batería}
\begin{itemize}
\item \textbf{Motor:} según el manual del operador de Lycoming (P/N 60297-12 o 60297-16) secciones 4 y 5 y la SI 1080 vigente.
\item \textbf{Silenciador:} abra la camisa de calefacción; sin grietas ni roturas, con atención al tubo de cola, uniones de los tubos verticales (\emph{risers}), soldaduras, abrazaderas, bridas y juntas; sin obstrucciones internas. \textbf{Prueba de fugas:} presurice con aire a baja presión y aplique agua jabonosa por fuera buscando burbujas.
\item Paneles de enfriamiento del motor (sobre todo el del enfriador de aceite y el de la manguera de enfriamiento del alternador) sin grietas ni sujetadores faltantes.
\item \textbf{Alternador y polea:} la polea debe ser de acero (verifique con un imán): no se aprueba la de aluminio. Cableado sin señales de arco; espárragos de salida, auxiliar y campo firmes; espaciador aislante sin daño. \textbf{Correa del alternador:} reemplazar si tiene grietas, dientes faltantes o delaminación; tensión según Lycoming SI 1129; alineación correcta. Retire la correa de repuesto de emergencia si está instalada. Manguera de enfriamiento sin obstrucciones ni agujeros.
\item Cortafuego vertical sin grietas, pandeos ni arrugas. Bastidores inferiores sin corrosión ni grietas; perno de cabeza perforada de la unión inferior izquierda bastidor--vientre con su torque (\S\,23-32) y alambre de frenado.
\item \textbf{Soportes del motor:} reemplace los de goma con deflexión significativa o grietas más que superficiales.
\item \textbf{Acelerador:} con colectivo abajo y empuñadura en ralentí, el balancín de la mariposa toca el tornillo de ajuste de ralentí; con colectivo a media carrera y acelerador abierto, toca el tope de apertura total. Holgura con el cortafuego vertical. Verifique el reglaje de la correlación (\S\,10.150) y del interruptor de acelerador a fondo (\S\,33-60).
\item \textbf{Mezcla:} el control lleva el brazo de tope a tope; abrazaderas de la funda del cable firmes (empuje y tire para comprobar que no resbale); el brazo gira libremente respecto del perno de retención del cable; el resorte de seguridad está instalado de modo que lleve el brazo a rica a fondo si se corta el cable.
\item Líneas de aceite y codos sin grietas ni abrasión y sin contacto con cables, amarras ni estructura; abrazaderas sanas; enfriador sin grietas ni fugas y con las aletas despejadas.
\item \textbf{Gascolator:} con la válvula en OFF, retire y limpie la cuba y el filtro; reinstale, abra la válvula, verifique que no haya fugas y frene el retén; drenaje que abra fácil, drene, cierre por resorte, selle y tenga franja de torque (\S\,23-31). Mangueras de combustible sin grietas ni fugas.
\item \textbf{Caja de aire:} disco de la válvula de drenaje libre (no aplica a los Mariner); compuerta alternativa con resorte libre y cierre completo. La calefacción abre y cierra por completo la válvula deslizante. Limpie el filtro B771-1 (\S\,6.435). Mangueras de entrada de aire, toma de calefacción y mangueras de calefacción sin obstrucciones, fugas, roces ni grietas.
\item Batería, cables y caja: electrolito (no selladas), prueba de capacidad o reemplazo según el fabricante, terminales sin grietas ni corrosión, tubo de drenaje despejado.
\end{itemize}

\subsection*{8. Carenado del cono de cola}
\begin{itemize}
\item Bastidor superior y unión con el cono de cola (inspección minuciosa de soldaduras).
\item \textbf{Abrazadera protectora del bastidor superior:} reemplácela si el desgaste la atravesó. Sin abrazadera, \textbf{el desgaste del tubo del bastidor por roce de las correas no debe exceder 0,0049~in}. La abrazadera puede instalarse en cualquier R22 (puede escotarse 0,12~in máximo en la guía para que su borde trasero quede al menos 0,05~in detrás del punto de contacto de la correa; separación de 0,01 a 0,05~in).
\item \textbf{Rodamiento superior del actuador:} juego axial máximo \textbf{0,060~in}; sin \emph{fretting} entre la pista interior y el eje del embrague; la pista no debe haber resbalado sobre el eje; Telatemp. Si hay un aumento inexplicable, inspeccione según \S\,22-41. Lubrique según el programa.
\item Montante de centrado lateral del embrague; placa flexible intermedia y horquillas (mismos criterios que la delantera); tubos del rotor de cola y balancín A331-1.
\item Correas en V según \S\,2.507 (sección siguiente).
\item \textbf{Eje del ventilador:} inspección visual de 360° de la parte expuesta en busca de grietas; pernos firmes y frenados.
\item \textbf{Rodamiento inferior del actuador:} la pista interior no debe haber resbalado sobre el eje; Telatemps; inspección \S\,22-42 si hay aumento inexplicable; lubricación según el programa.
\item \textbf{Soportes A185 del rodamiento inferior:} sin \emph{fretting} ni holgura en soportes y remaches (si hay holgura, repare según R22 SL-74A). Aplique una franja horizontal de sellador de torque sobre los dos remaches exteriores del soporte inferior hacia la carcasa inferior.
\item Carcasa de fibra de vidrio sin daños en las paletas de la carcasa superior derecha; agujero de drenaje despejado. \textbf{Labios metálicos de entrada:} separación de 0,020 a 0,100~in con la entrada del ventilador.
\item \textbf{Ventilador:} sin grietas ni corrosión; bordes de ataque de las paletas sin daños; \textbf{pasador elástico alineado con las marcas de deslizamiento} (si no, retire el ventilador e inspeccione cubo y eje).
\item \textbf{Actuador A051-1:} con el embrague acoplándose, presione la palanca del interruptor de límite de extensión: el motorreductor debe detenerse y reanudar al soltarla. Con un espejo, observe los resortes de columna al final del ciclo: deben saltar hacia afuera simultáneamente. No exceder el límite de extensión acoplada (figura 7-15 del MM). Holgura mínima de 0,030~in entre el motorreductor y el bastidor superior.
\item \textbf{Interruptores de los resortes del actuador:} método 1, con el enchufe de prueba MT558-1 en el conector ``Test'' del arnés (puentes 1-2 para el interruptor del cable 98 y 2-3 para el del cable 91): el motorreductor debe detenerse al saltar los resortes; si no se detiene, el interruptor falló cerrado: retire el MT558-1 de inmediato. Método 2, sin enchufe de prueba: hunda los resortes de un lado hasta que salten hacia adentro y manténgalos al menos 1~s; el motor no debe arrancar (si arranca, déjelo unos 2~s, suelte y repita: si vuelve a arrancar, falla el interruptor opuesto); desacople, vuelva a acoplar y repita del otro lado. Reemplace el interruptor que falle (\S\,7.551) antes de volar.
\item \textbf{Polea inferior:} superficies lisas; elimine la corrosión superficial y retoque con imprimante aprobado (de preferencia cromato); reemplácela si tiene picaduras, descascarado del metalizado, asperezas o crestas filosas.
\item Verifique la altura del motor (\S\,6.130), el ángulo del eje del embrague (\S\,7.240; reemplace las correas si es necesario), la alineación de poleas (\S\,7.230) y el reglaje de la correlación del acelerador (\S\,10.150).
\end{itemize}

\subsection*{9. Tapas de inspección del cono de cola}
\begin{itemize}
\item Verifique el autobloqueo de las tuercas remachadas al retirar las tapas.
\item \textbf{Eje del rotor de cola:} sin grietas, dobleces ni contacto con el interior del cono; corrosión dentro de los límites (figura 7-12C del MM); descentramiento según \S\,7.340.
\item Tubo del rotor de cola y sus terminales; los vástagos de los tornillos del protector del rotor de cola no deben tocar el tubo.
\item \textbf{Amortiguador:} firme; \textbf{desgaste máximo de 0,040~in} en la carcasa del rodamiento.
\item Interior del cono sin remaches sueltos ni corrosión; bujes del tubo sin desgaste excesivo; sin grietas, sobre todo en la unión del amortiguador.
\item Antena de látigo inferior (retire los cuatro tornillos de la placa y revise grietas y corrosión) y antenas superiores.
\item Exterior del cono (\S\,4.310): mellas, rayas, abolladuras, grietas, corrosión, \emph{fretting} o remaches sueltos; drenajes al frente de cada bahía (excepto la delantera) despejados; grietas junto a los montajes de antena.
\item Estabilizador horizontal B375-2, protector visual del rotor de cola y luz estroboscópica.
\end{itemize}

\subsection*{10. Cubierta plástica}
\begin{itemize}
\item Límpiela y verifique el autobloqueo de las tuercas al reinstalarla.
\item Empenaje sin grietas ni sujetadores sueltos; drenajes del estabilizador vertical inferior y del patín de cola despejados. Luz de navegación trasera.
\item Placa flexible trasera (mismos criterios).
\item \textbf{Caja de cola:} sellos, nivel de aceite sin fugas, alambre de frenado, eje de salida limpio. Cada 500~h o anualmente, limpie el detector de virutas (\S\,22-21); drene y lave cada 500~h máximo (\S\,22-23).
\item \textbf{Control de paso del rotor de cola y balancín trasero:} se permite \emph{una} grieta radial en la bola del rodamiento esférico inferior; \textbf{juego rotacional máximo de 0,25~in} medido en la cabeza del perno de la varilla de paso; el conjunto debe deslizar libremente sobre el eje de salida; sellos sin fugas.
\item \textbf{Varillas de paso:} terminales y rodamientos según \S\,2.120; pueden invertirse (extremo exterior hacia adentro) para prolongar su vida; puede instalarse un anillo O A215-012 opcional en el espaciador A115-1 bajo la cabeza del perno. Reinstale los contrapesos de cuerda en sus puntos.
\item \textbf{Palas del rotor de cola} (R22 SB-120; límites en \S\,28-30): sin pandeo; sin holgura entre las pistas exteriores de los rodamientos de cambio de paso y los herrajes de raíz; drenajes de punta despejados.
\item \textbf{Cubo del rotor de cola:} sin mellas, rayas, excavaciones, abolladuras, grietas ni corrosión; rodamientos elastoméricos según \S\,2.125; al balancear el cubo, el perno, los espaciadores y las tuercas no deben girar; balanceo suave; pernos de pala firmes.
\end{itemize}

\begin{advertencia}
Los controles de paso A031-1 rev.~J o anteriores requieren inspección repetitiva según el R22 SB-90A cada 300 horas o anualmente, lo que ocurra primero. Puede producirse daño estructural si se aplica aire comprimido a los agujeros de drenaje de las puntas de las palas.
\end{advertencia}

\subsection*{11. Carenado del mástil}
\begin{itemize}
\item Carenado sin grietas, abolladuras ni corrosión; sujetador A665-1 sin ceder ni agrietarse (puede ceder por apretar demasiado los tornillos de sus tuercas remachadas). Costillas superior e inferior firmes (sobre todo la abrazadera inferior).
\item \textbf{Tres tubos verticales} (dos A121-7 y un A121-5): sin grietas en los extremos superiores; con colectivo arriba y cíclico atrás, sin desgaste de los A121-7 ni de sus fundas donde rozan en la guía.
\item Tubo pitot y su línea: agujero de drenaje bajo el codo detrás del pitot despejado. Respiraderos de los tanques sin dobleces ni obstrucciones.
\item \textbf{Tijeras superiores del plato oscilante:} mida el juego (con cojinetes de muñón C648-2: 0,001 a 0,010~in de juego axial total por pivote; con rodamientos plásticos A139 o esféricos B115: 0,060~in total). Con alguien subiendo y bajando el colectivo, el perno, los muñones (o las bolas y espaciadores) y el brazo deben girar juntos en cada pivote.
\item Tijeras inferiores. \textbf{Tubo deslizante:} sin grietas en los agujeros de remaches ni corrosión en la base; sin daño ni desgaste del anodizado (un leve ``cuadriculado'' bajo lupa de 10X es normal).
\item Retire la amarra inferior de la funda del plato oscilante y, con espejo, revise entre el eje del rotor y el interior del tubo deslizante: sin corrosión ni residuos.
\item \textbf{Plato oscilante:} juego radial máximo de \textbf{0,020~in} entre la bola y el tubo deslizante; al girar el rotor a mano, sin rodamientos ásperos o secos. \textbf{Suplementado:} observando la bola desde abajo con alguien moviendo lentamente el colectivo, la bola debe moverse de inmediato con el plato al invertir el sentido; si no, hay juego axial y debe suplementarse (\S\,8.416).
\item Reinstale la amarra de la funda. Eje del rotor principal sin corrosión.
\end{itemize}

\subsection*{12. Cubo y pernos de bisagra}
\begin{itemize}
\item Cubo sin mellas, rayas, excavaciones ni corrosión; sin residuo café o negro (desgaste de rodamientos). Si se sospecha desbalance, revise la fricción de las bisagras de balanceo y conicidad (\S\,26-32).
\item \textbf{Pernos de bisagra (tres): la corrosión está prohibida}; chavetas instaladas; cabezas y tuercas con franja de torque hasta las arandelas de empuje.
\item Varillas de paso y terminales: centrado, contratuercas apretadas y alambre de frenado correcto.
\end{itemize}

\subsection*{13. Palas del rotor principal}
\begin{itemize}
\item Fundas sin daños ni fugas, bien posicionadas y con holgura respecto del cubo en todo el recorrido.
\item Espigas y herrajes de raíz según \S\,26-43; AD 88-26-01 (solo espigas A158-1).
\item Cuidado del borde de ataque (\S\,26-53) y mantenimiento de las puntas (\S\,26-60).
\item Revestimientos y dobladores: rayas y corrosión (\S\,26-41); abolladuras y deformaciones (\S\,26-42); huecos (\S\,26-44). Encere con cera tipo carnaúba. Drenajes de punta despejados.
\end{itemize}

\subsection*{14. Cabina}
\begin{itemize}
\item Sin equipos sueltos que puedan interferir con los controles.
\item \textbf{Cinturones:} sin deshilachado, costuras rotas, agujeros ni podredumbre; sin decoloración significativa por UV; carretes inerciales y hebillas funcionando. (Los arneses instalados de fábrica no requieren etiqueta TSO.)
\item \textbf{Ventanas:} sin daño UV significativo (amarillamiento). Defectos aceptables: una mella de no más de 0,010~in de profundidad y no mayor de 0,25 × 0,50~in por pie cuadrado; rayas de no más de 0,010~in de profundidad y 5~in de largo; manchas que se eliminen con un pulido ligero. Para grietas y cuarteado junto a los bordes del parabrisas, vea \S\,25-20.
\item \textbf{Hilo de guiñada: mínimo 3~in de largo a cada lado.}
\item Exterior sin grietas, corrosión ni remaches sueltos (pintura agrietada o residuo negro alrededor de las cabezas). Luces de navegación roja (izquierda) y verde (derecha).
\item Puertas: sin grietas junto a bisagras y pestillos; ventilaciones funcionando; pasadores de bisagra con chaveta; tornillos de bisagras apretados; pestillos y cerraduras funcionando.
\end{itemize}

\subsection*{15. Tren de aterrizaje}
\begin{itemize}
\item \textbf{Tubos de patín y zapatas:} zona plana por desgaste bajo el tubo de no más de 0,5~in de ancho; rayas longitudinales de no más de 0,030~in de profundidad; rayas circunferenciales de no más de 0,030~in de profundidad y menos de 0,50~in de largo; drenajes abiertos (no aplica a flotadores); tapas de lluvia firmes; \textbf{espesor mínimo de zapata 0,05~in}.
\item Montantes o soldaduras sin grietas ni corrosión. Si el helicóptero hace habitualmente aterrizajes corridos, verifique el torque de los pernos montante--patín.
\item \textbf{Tubos transversales:} altura mínima del patín de cola (\S\,5.210); sin corrosión (si una tapa de lluvia está suelta, revise la corrosión interna); sin grietas ni \emph{fretting} en las uniones con los montantes. Un residuo gris oscuro indica \emph{fretting}: marque el tubo a 1~in hacia adentro de la unión, retírelo y haga inspección con penetrantes fluorescentes del extremo por fuera de la marca.
\item Puntos de unión del tren sin sujetadores sueltos, grietas, pandeo ni \emph{fretting}; soportes sin estampados sueltos ni rodamientos gastados.
\item Flotadores: pruebas de la válvula de alivio y de mamparos (\S\,2.580); presión según el POH. Tubos de patín de flotadores: retire los tapones de drenaje traseros y, con las ruedas instaladas, tire hacia abajo el cono de cola para drenar la humedad.
\end{itemize}

\subsection*{16 a 19. Partes de vida limitada, AD, SB, documentos y registros}
\begin{itemize}
\item Reemplace las partes de vida limitada según los registros y \S\,3.300; verifique que los componentes instalados tengan tiempo suficiente para la operación prevista.
\item Verifique el cumplimiento de las AD y los SB aplicables a célula, motor y accesorios (algunos exigen inspecciones recurrentes a menos de 100~h o anuales).
\item A bordo, legibles y vigentes: certificado de aeronavegabilidad, matrícula, licencia de estación de radio si corresponde, POH y lista de equipos/peso y balance. Letreros instalados y legibles (sección 2 del POH).
\item Retire herramientas, ferretería suelta, paños y objetos extraños; cierre todas las tapas. En los Mariner, selle el perímetro de la tapa de la caja de aire con cinta de aluminio.
\item Registre el trabajo con descripción (o referencia a datos aceptables), fecha, tiempo total, firma, tipo y número de licencia de quien aprueba el retorno al servicio. La lista incluye casillas para anotar la alineación de poleas izquierda y derecha, la altura del motor, el ángulo del embrague y el descentramiento del eje del rotor de cola.
\end{itemize}

\section{Inspecciones especiales}

\subsection{Correas en V}

\begin{table}[htbp]
\centering\small
\caption{Discrepancias de las correas en V, causas y soluciones (MM \S\,2.507).}
\begin{tabularx}{\linewidth}{@{}L{2.2cm}YY@{}}
\toprule
Síntoma & Causa & Solución \\
\midrule
Correas agrietadas & Correa vieja & Reemplazar correas (\S\,7.280) \\
 & Correa que patina: calor y endurecimiento progresivo de la base & Reemplazar correas y el actuador (\S\,7.500) \\
Parte superior de la banda de unión deshilachada o dañada & Obstrucción que interfiere con la correa & Eliminar la obstrucción y reemplazar correas \\
Banda de unión ampollada o perforada & Tierra o arena acumulada entre correas y polea & Reemplazar correas; lavar periódicamente el interior de las nuevas \\
 & Correas gastadas & Reemplazar correas \\
Correa cortada por dentro & Forzada sobre el reborde de la polea al instalarla sin holgura & Reemplazar correas \\
 & Material extraño cayó en la transmisión & Retirarlo y reemplazar correas \\
Correa montada fuera de la ranura & Forzada sobre el reborde al instalarla sin holgura & Reemplazar correas \\
\bottomrule
\end{tabularx}
\end{table}

\subsection{Patrón de desgaste en la polea inferior}

Observe los patrones de desgaste en el imprimante de las cuatro ranuras de la polea inferior: los ocho patrones (ambos flancos de las cuatro ranuras) deben ser muy similares y uniformes en toda la circunferencia. Si difieren notoriamente, mida el ancho del más angosto y del más ancho: \textbf{si la relación excede 3:1, reemplace las correas}. Profundidad de ranura: 0,60~in en las poleas inferiores ``de ranura profunda'' (rev.~S y anteriores) y 0,43~in en las ``de ranura baja'' (rev.~T y posteriores); ancho de la polea 0,925 y 0,735~in respectivamente.

\begin{precaucion}
RHC no recomienda usar poleas superiores A170-1 o -2 rev.~N o anteriores con poleas inferiores A493-1 rev.~T o posteriores: pueden causar vibración excesiva de las correas.
\end{precaucion}

\subsection{Golpe del patín de cola}

\textbf{A. Raspadura o leve flexión del patín.} Inspeccione visualmente las palas del rotor de cola (si hay daño, \S\,2.520); el estabilizador vertical inferior (sin pandeo, grietas ni remaches sueltos en las uniones con el patín y el estabilizador horizontal); las uniones del estabilizador horizontal al cono, y el cono y sus uniones con el bastidor superior.

\textbf{B. Pandeo o rotura del patín o del estabilizador vertical inferior.} Además de A: descentramiento del eje del rotor de cola (\S\,7.340); retire el cono de cola y verifique que los agujeros de unión ovalados no superen \textbf{0,265~in} en su parte más ancha; decape 2~in alrededor de los puntos de unión en el bastidor superior e inspeccione con penetrantes fluorescentes; retire el conjunto estabilizador (agujeros $\le$ 0,265~in, sin remaches sueltos ni pandeo) y haga penetrantes de la fundición del cono; separe el estabilizador vertical inferior del horizontal: el pandeo o las grietas exigen reemplazo; los remaches sueltos pueden reemplazarse; agujeros ovalados $\le$ \textbf{0,200~in}.

Además, ante daño en los estabilizadores A042-1, A043-1 o A044-1 (abolladura, rotura o grieta del metal) o tras un golpe del patín, decape las superficies indicadas del mamparo A148 del cono de cola, inspeccione con penetrantes (reemplace el cono si hay grieta), aplique conversión química, imprimante epóxico y pintura (no en las superficies de unión de la caja de cola).

\subsection{Golpe del rotor de cola}

\textbf{A. Contacto con piedras pequeñas, pasto alto u otro objeto pequeño en el aire.} Inspeccione las palas (\S\,28-30) y el descentramiento del eje: \textbf{si excede 0,025~in en cualquier punto, el eje debe reemplazarse o repararse}.

\textbf{B. Detención súbita por contacto con el suelo u objeto sólido} (flexión o corte de palas). Además de A: decape al menos 2~in alrededor de los agujeros delanteros y la soldadura trasera del eje e inspeccione con penetrantes (grieta = reemplazo); deformaciones o abolladuras del eje = reemplazo (las mellas y rayas pueden pulirse, \S\,7.360); envíe el rotor de cola y la caja de cola a RHC para overhaul; reemplace las placas flexibles trasera e intermedia; verifique que las palas no hayan golpeado el cono ni el empenaje, e inspeccione el sistema rotor principal.

\subsection{Golpe del rotor principal}

\textbf{A. Contacto con objetos pequeños en el aire.} Inspeccione las palas (\S\,26-40) y verifique que el borde de fuga no esté pandeado ni doblado.

\begin{advertencia}
Cualquier pandeo de una pala o flexión del larguero se considera detención súbita y exige cumplir la parte B.
\end{advertencia}

\textbf{B. Detención súbita.} Verifique el descentramiento del eje del rotor de cola; retire y envíe a un taller autorizado por RHC: palas y espigas A005-2, cubo y rodamientos A005-4, plato oscilante A017-2, embrague A018-1 o -2, muñones A106, arandelas de empuje A152-1 y horquillas A907-4/-5 y A908-4 o A006-5. Reemplace la placa flexible delantera A947-1, envíe la caja principal a RHC para overhaul e inspeccione el motor según las instrucciones de detención súbita de su fabricante.

\subsection{Sobrevelocidad del rotor o del motor}

\textbf{A. Rotor entre 110 y 116\,\%} (si fue con potencia, vea también C):
\begin{pasos}
\item Verifique el balance dinámico de los rotores principal y de cola y compárelo con el previo. \textbf{Un cambio mayor que 0,3~ips exige la inspección B.}
\item Retire las palas, drene las carcasas de los rodamientos de paso, retire las abrazaderas exteriores de las fundas y gire las espigas para verificar que no haya \emph{brinelling} (identaciones) en los rodamientos. (Tienen alta precarga; una leve aspereza es normal; si es evidente, envíe el conjunto a un taller autorizado.)
\item Inspeccione visualmente las palas principales y de cola.
\item Verifique el descentramiento del eje del rotor de cola.
\end{pasos}

\textbf{B. Rotor a 116\,\% o más:} pasos 2 a 4 de A; envíe palas y cubo a un taller autorizado para ensayos no destructivos de espigas y cubo; reemplace pernos, arandelas de empuje y muñones agrietados o doblados e inspeccione con partículas magnéticas los que no se reemplacen.

\textbf{C. Sobrevelocidad del motor.} El porcentaje de sobrevelocidad se calcula a partir de la indicación del tacómetro:
\[
\%\ \text{sobrevelocidad} = \frac{\text{indicación (\%)} \times 2550}{2700} - 100
\]
(104\,\% indicado = 2652 RPM reales; el motor está clasificado a 2700 RPM). Para la inspección, vea la Lycoming SB 369 vigente.

\subsection{Aterrizaje duro}

\textbf{A. Deformación de los tubos transversales sin cargas laterales} (las cargas laterales se reconocen por pandeo o flexión del bastidor de tubos o de la chapa):
\begin{pasos}
\item Inspeccione los bastidores de tubos (sin deformación ni grietas), con atención a los montantes verticales traseros del bastidor inferior.
\item Descentramiento del eje del rotor de cola (\S\,7.340).
\item Palas del rotor principal (\S\,26-40).
\item Con el helicóptero en suelo nivelado, mida del patín de cola (o de la base del estabilizador del flotador) al suelo. \textbf{Reemplace el tubo transversal trasero} si la medida es menor que: 24~in (Standard o HP); 34~in (Alpha, Beta, Beta II, Mariner o Mariner II, desde el patín); 37~in (Mariner o Mariner II, desde el estabilizador del flotador). \textbf{Reemplace el tubo transversal delantero} si toca su cubierta.
\item Verifique y ajuste la alineación de las poleas (\S\,7.230).
\end{pasos}

\textbf{B. Aterrizajes duros con posibles cargas laterales:} todo lo de A, más: puntos de unión del tren al fuselaje y cortafuegos vertical y horizontal sin pandeo ni grietas (atención a las uniones de los montantes de soporte del motor al cortafuego vertical); retire el cono de cola y haga penetrantes fluorescentes del bastidor superior en los puntos de unión del cono (decapando 2~in alrededor); inspeccione el cono (sin pandeo ni remaches sueltos) y los codos del tren (sin deformación, mellas ni grietas).

\subsection{Corrosión del bastidor de tubos de acero}

\begin{pasos}
\item Para áreas grandes, retire el bastidor y decape para determinar la extensión. Si excede los límites, reemplácelo.
\item Mida y registre el diámetro del tubo en la zona corroída.
\item Pula la corrosión con lija al agua de grano 320 o más fino, vuelva a medir y verifique que la profundidad pulida esté dentro de límites.
\item Limpie, imprima y pinte (\S\,23-71); reinstale el bastidor.
\end{pasos}

\begin{checklist}
\chk{Corrosión superficial leve: profundidad máxima tras el retrabajo}{0,002 in}
\chk{Picaduras: profundidad máxima tras el retrabajo}{0,004 in}
\chk{Picaduras: diámetro máximo}{0,10 in}
\chk{Separación mínima entre zonas con picaduras}{0,50 in}
\end{checklist}

\subsection{Inspección y prueba de flotadores}

Haga la prueba de la válvula de alivio y la de mamparos en cada inspección de 100 horas o anual. Los flotadores se reparan o reemplazan según su estado o si fallan una prueba.

\begin{precaucion}
Abra la válvula de inflado antes de aplicar presión (no aplique presión contra una válvula cerrada). Antes de la prueba, inspeccione el flotador desinflado y con cada compartimiento inflado individualmente a 0,75--1,0~psig: sin cortes, roturas, perforaciones, abrasión, delaminación (ampollas) ni separación de costuras. No exceda 7~psig en los compartimientos.
\end{precaucion}

\textbf{Válvula de alivio.} Infle los compartimientos a 5,0~psig con aire seco, filtrado y sin aceite (con múltiple o de a 0,5~psig por cámara). Llene con agua limpia las aberturas de las válvulas, aumente lentamente la presión y registre la presión de apertura (burbujas); luego reduzca hasta que cesen y registre la de cierre. Reemplace la válvula si la apertura no está entre \textbf{5,5 y 6,5~psig} o el cierre es menor que \textbf{5,5~psig}.

\textbf{Mamparos.} Desinfle todo e infle la cámara delantera y la segunda desde atrás a 3~psig. Si tras 1 hora a temperatura constante la presión es menor que \textbf{2,75~psig}, repare o reemplace. Repita con la cámara trasera y la segunda desde adelante.

\section{Inspección de 2200 horas / 12 años}

\begin{nota}
Los kits de inspección de 2200 horas son el KI-2202 (S/N 2570 y anteriores) y el KI-2203 (S/N 2571 y posteriores). La inspección de 12 años solo se requiere en helicópteros con 12 años en servicio y menos de 2200 horas desde nuevos, desde la última inspección de 2200 horas o desde la última de 12 años.
\end{nota}

\begin{pasos}
\item Reemplace las partes de vida limitada (o sus conjuntos superiores) según los registros y \S\,3.300.
\item Efectúe el mantenimiento adicional de componentes (\S\,1.102; capítulo~\ref{cap:mm-vida-limitada}).
\item Retire el motor, los paneles de enfriamiento, el ventilador, la carcasa, la inducción, la calefacción del carburador, las líneas de aceite y todas las mangueras. Haga overhaul o inspección de magnetos, alternador y carburador según sus fabricantes.
\item Retire los estabilizadores (\S\,4.320) y reemplácelos si corresponde; si no, inspeccione (sin grietas, corrosión, remaches sueltos, abolladuras ni deformación) y haga penetrantes de las zonas sospechosas.
\item Retire el tren (\S\,5.100) y reemplácelo si corresponde; si no, desármelo por completo, decápelo, inspecciónelo (partículas magnéticas y penetrantes), imprima, pinte y rearme.
\item Retire los bastidores de tubos de acero (\S\,4.200) y reemplácelos si corresponde; si no, decape, inspeccione con partículas magnéticas, imprima y pinte.
\item Retire el bastón del cíclico (\S\,8.110) y su arnés; decape por chorro de medio seco, inspeccione con partículas magnéticas, imprima, pinte y reinstale.
\item Retire los tubos A121 (capítulo 8 del MM), decape, inspeccione con penetrantes, imprima, pinte y reinstale.
\item Inspeccione con partículas magnéticas los herrajes de montaje A650 y el perno F650-3 de la caja principal y el eje intermedio A337-1.
\item Inspeccione el cableado de la célula: sin corrosión ni deterioro del aislamiento; cada cable en su disyuntor correcto.
\item Retire las vejigas de combustible principal y auxiliar (\S\S\,12.110 y 12.120): sin grietas ni delaminación del caucho ni de sus uniones con el metal; insertos roscados y acoples sin daño; presurice a \textbf{1~psi máximo} y busque fugas con agua jabonosa; revise las estructuras de los tanques por evidencia de fugas; reinstale.
\item Alineación de poleas (\S\,7.230) y suplementado de la placa flexible intermedia (\S\,7.330).
\item Reglaje de controles y ángulos de pala del rotor principal (\S\S\,10.110 y 10.120) y del rotor de cola (\S\S\,10.130 y 10.140).
\item Inspección de 100 horas/anual (\S\,2.400).
\item Pese el helicóptero (\S\,18-20).
\item Balancee el rotor de cola (\S\,10.240); \emph{tracking} y balanceo del rotor principal (\S\,10.230); chequeos en tierra, prueba de motor y vuelo (\S\S\,2.205, 2.210 y 2.220). (La operación prolongada a baja potencia con aros de pistón nuevos puede impedir su asentamiento.)
\item Drene y lave las cajas principal y de cola (\S\S\,22-13 y 22-23).
\item Haga las anotaciones de mantenimiento correspondientes.
\end{pasos}

\subsection{Inspección y reparación del cubo del rotor principal}

\begin{pasos}
\item Retire los rodamientos del cubo (\S\,26-21 pasos 1 y 2).
\item Decape el cubo por chorro de medio plástico o bicarbonato.
\item Pula con lija de óxido de aluminio de grano 320 o más fino, con radio de transición mínimo de 0,25~in:
\begin{itemize}
\item Alojamientos de rodamientos: sin rayas ni daños evidentes. \emph{Fretting} o corrosión: máximo \textbf{0,002~in} en las superficies interior y exterior dentro de 1,750~in de los alojamientos (1,875~in en los de bisagra de conicidad del cubo B370-1); máximo \textbf{0,0005~in} dentro de los alojamientos. Diámetro máximo del alojamiento: \textbf{1,1870~in} (1,3120~in en los de conicidad del B370-1).
\item Resto del cubo: mellas, rayas, excavaciones o corrosión hasta \textbf{0,010~in} dentro de 0,30~in de los bordes interior y exterior, y hasta \textbf{0,060~in} en las demás áreas.
\end{itemize}
\item Inspeccione con penetrantes fluorescentes (\S\,23-42).
\item Imprima (\S\,23-60), instale los rodamientos (\S\,26-21 pasos 4 y 5), enmascare los rodamientos y aplique la capa final; retire el enmascarado una vez curada.
\end{pasos}

\subsection{Mantenimiento de 12 años de la caja de cola B021-1}

Drene y lave la caja (\S\,22-23) y retírela (\S\,7.410). Retire la horquilla B546-1 del engranaje de entrada B545 e inspecciónela; inspeccione las estrías del engranaje (si están dañadas o con escalones de desgaste, envíe la caja a RHC); reemplace el sello del eje de entrada (\S\,7.470). Retire el visor, limpie el vidrio e inspeccione el juego de engranajes (corrosión o daño: envíe a RHC); reinstale el visor con su torque especial. Reemplace el sello del eje de salida (\S\,7.460). Instale la caja, llene hasta el centro del visor con el aceite correcto, gire el rotor varias vueltas y tire hacia abajo de la caja varias veces, reajuste el nivel y apriete el tapón con su torque especial. Conecte la luz de cola y el detector de virutas y pruebe ambos (\textsf{TR CHIP}). Instale el rotor de cola (\S\,28-10).

\subsection{Mantenimiento de 12 años del ventilador B174-1}

Retire el ventilador (\S\,6.210) y marque todas sus piezas, incluida la ferretería de balanceo, para rearmarlo idéntico. Retire el cono A187-4, el cubo A186-2 y los espaciadores A178-2. Inspeccione todas las piezas (sin daños ni agujeros ovalados) y reemplace las necesarias. Rearme con imprimante aprobado húmedo (de preferencia cromato) entre todas las superficies de apriete y, con el imprimante aún húmedo, aplique torque especial: \textbf{pernos NAS6603, 70~in-lb; pernos NAS6605, 300~in-lb} (12 pernos NAS6603-6 con arandelas NAS1149F0332P bajo la cabeza, NAS1149F0316P entre cono y ventilador, cuatro NAS1149D0363J bajo cada tuerca D210-3; 6 pernos NAS6605-12, -13 o -14 ---los necesarios para dejar 2 a 4 hilos expuestos--- con arandelas NAS1149C0516R entre cubo y cono y NAS1149F0532P bajo la tuerca D210-5).

\subsection{Balance estático de las palas del rotor principal}

\begin{precaucion}
Acolche suficientemente las superficies que tocan los revestimientos y bordes de fuga para no dañar las palas.
\end{precaucion}

\begin{pasos}
\item Arme las palas con espigas en el cubo eligiendo muñones de bisagra de conicidad A106 y suplementos A117 para lograr una \textbf{holgura axial total de 0,002 a 0,006~in}; instale el muñón del mismo largo (número de guion) en el lado del borde de ataque de ambas bisagras.
\item Instale las tuercas en seco y apriételas hasta asentar muñones y arandelas; mida las separaciones entre arandelas y caras de los rodamientos y verifique la holgura en cada bisagra.
\item Instale el útil de balance estático MT009-10 en el cubo, ice el rotor y apóyelo en la base MT009-9.
\item Verifique que los ángulos de paso de ambas palas sean aproximadamente iguales (si no, los cuernos de paso están mal instalados).
\item Nivele en cuerda (paralelo a las bisagras) con un nivel de burbuja sobre el cubo, ajustando los tornillos NAS1351-3-10P.
\item Con un micrómetro de profundidad o pie de metro, mida la distancia entre la placa extrema MT009-2 y el eje MT009-8 a ambos lados del cubo; marque con ``N'' el lado de menor distancia (lado de la tuerca del perno de balanceo).
\item Nivele en envergadura (perpendicular a las bisagras) colocando la tapa de punta, sus tornillos y dos tornillos A722-4 lo más cerca posible de la punta y agregando tiras de balance A298 y arandelas NAS1149F0332P o NAS1149F0363P. \textbf{El balance final debe quedar dentro de una arandela NAS1149F0332P.}
\item Aplique antiagarrotante A257-9 en las roscas de los tornillos de los pesos de punta y de la tapa de punta y apriételos con torque especial de \textbf{40~in-lb}; vuelva a verificar el balance.
\item Marque el conjunto con ``X'' y ``O'' de color en el cubo, las raíces de las palas y los pernos de conicidad, como referencia de armado; desarme.
\end{pasos}

\subsection{Inspección de las varillas de paso A258-5}

Marque cada varilla y su oreja del plato oscilante con un color y registre sus longitudes. Retírelas y desármelas sin mezclar piezas. Evalúe los terminales (\S\,2.122); inspeccione las piezas con lupa de 10X y el interior de los cilindros con boroscopio; reemplace lo dañado. Rearme a las longitudes registradas; apriete contratuerca y \emph{palnut} (\S\,23-32) y aplique el torque especial a la contratuerca autoblocante (\S\,23-33). Instale, apriete y marque con franja de torque.

\section{Inspecciones especiales adicionales}

\subsection{Ceniza volcánica}

Debe evitarse el vuelo en condiciones de ceniza volcánica visible (``nube de ceniza''). Si se operó en ellas:

\begin{pasos}
\item Según Lycoming SI 1530, con equipo de protección, retire la mayor cantidad de residuo con aspiradora y luego aire comprimido (no cerca de los drenajes de las palas).
\item Limpie, lave y enjuague a fondo el helicóptero (sección 8 del POH), incluida la cara interior de las correas.
\item Retire las tapas de punta de las palas y limpie las puntas.
\item Con lupa de 10X, inspeccione la línea de unión revestimiento--larguero expuesta: \textbf{la pala no es aeronavegable si hay cualquier hueco, incluso ``agujeros de alfiler''}. Repinte según sea necesario.
\item Inspeccione las poleas; reemplace las que tengan picaduras, descascarado, desgaste del recubrimiento, asperezas o crestas filosas, y las correas si alguna polea tiene crestas filosas.
\item Desconecte la correa del alternador y gire su polea a mano: rodamientos y escobillas deben girar suavemente (la ceniza puede entrar por el aire de enfriamiento sin filtrar). Revise las poleas del alternador y del soporte de la corona por escalones de desgaste (reemplace correa y poleas si los hay) y ajuste la tensión (SI 1129).
\item Limpie el interior de la caja de aire e inspeccione el filtro. Revise la inducción aguas abajo del filtro (un guante blanco limpio ayuda); si hay ceniza: limpie la inducción, desarme e inspeccione el carburador (overhaul si es necesario), haga la prueba de compresión (SI 1191), la evaluación de válvulas de escape (SB 388) y revise las bujías.
\item Retire el tapón de inspección del engranaje distribuidor de cada magneto; si hay ceniza dentro, haga overhaul (sus respiraderos no tienen filtro).
\item Revise el aceite: si huele mal, está opaco (o no claramente café) o hay partículas en la varilla, cambie aceite y filtro, inspeccione la malla de succión y el filtro usado y haga la prueba de compresión.
\end{pasos}

\subsection{Impacto de rayo}

Son extremadamente raros en helicópteros VFR. Si ocurre, RHC recomienda una inspección de 100 horas y seguir la Lycoming SB 401. El alto voltaje bien conducido por la estructura se disipa con poco daño; mal conducido puede causar calor excesivo que hornee, queme, carbonice o funda materiales. Sobre 300~°F algunos materiales pierden resistencia aunque no muestren daño visible. Inspeccione a fondo las palas, el tren, la transmisión, la célula y los controles buscando arcos, quemaduras, picaduras o grietas, sobre todo en terminales de varilla, muñones y otros puntos de la ruta conductora. Si hay daño evidente, consulte a RHC con documentación detallada antes de aprobar el retorno al servicio.
